\section{Алгоритм поиска оптимального набора параметров}

Поиск оптимального набора параметров требуется для каждой частоты инжектируемого сигнала из исходного набора частот ${f^{(k)}_{z}}$, ${k \in \left[1;K \right]}$. В результате определяются такие $K$-наборов параметров, при котором частотная характеристика измеряется алгоритмом \ref{sec:fr_optim} за минимальное время с точностью воспроизведения частоты ${f^{(k)}_{z}}$ близкой к заданной ${\Delta\epsilon^{*}_{f}}$.

\textbf{Входные данные:}
\begin{itemize}
    \item Частота инжектируемого сигнала: ${f_z}$.
    \item Максимальное абсолютное отклонение от заданной частоты ${f_z}$ инжектируемого сигнала: ${\Delta\epsilon^{*}_{f}}$.
    \item Частота выполнения операций в системе: ${f_d}$.
    \item Максимальное количество точек дискретных последовательностей ${x[n]}$, ${y[n]}$ на интервале измерения: ${N_\text{max}}$.
    \item Минимальное количество точек дискретных последовательностей ${x[n]}$, ${y[n]}$ на интервале измерения: ${N_\text{min}}$.
\end{itemize}

\textbf{Выходные данные:}
\begin{itemize}
    \item Количество точек дискретных последовательностей ${x[n]}$, ${y[n]}$ на интервале измерения: ${N}$.
    \item Коэффициент децимации: ${M}$.
    \item Коэффициент кратности длительности измерительного окна ${T_h}$ к периоду инжектируемого сигнала ${T_z}$: ${L}$.
\end{itemize}

\textbf{Последовательность действий:}
\begin{enumerate}
    \item Вычисляется коэффициент соотношения между периодом инжектируемого сигнала и периодом выполнения операций в системе:
        \begin{equation*}
            R = \frac{f_{d}}{f_{z}}.
        \end{equation*}
    \item Вычисляется максимально возможное значение коэффициента децимации:
        \begin{equation*}
            M_\text{max} = \left\lfloor \frac{R}{2} \right\rfloor.
        \end{equation*}
    \item В цикле ${M \in \left[ 1; M_\text{max} \right]}$:
    \begin{enumerate}
        \item Вычисляются всевозможные значения
            \begin{equation*}
                L = \frac{M \cdot N}{R}
            \end{equation*}
        для ${N \in \left[ N_\text{min}; N_\text{max} \right]}$, с округлением $L$ до 1 при ${L<1}$.
        \item Значения $L$ используются для вычисления всевозможных приближений к начальному значению коэффициента ${R}$:
            \begin{equation*}
                R^{'} = \frac{M \cdot N}{\left\lfloor L \right\rfloor}.
            \end{equation*}
        \item Значения ${R^{'}}$ используются для вычисления всевозможных приближений к исходному значению частоты ${f_{z}}$:
            \begin{equation*}
                f^{'}_{z} = \frac{f_d}{R^{'}},
            \end{equation*}
            и соответствующих отклонений
            \begin{equation*}
                \Delta\epsilon_{f} = \left| f^{'}_{z} - f_{z} \right|.
            \end{equation*}
        \item Из набора значений ${f^{'}_{z}}$ исключаются те, для которых ${\Delta\epsilon_{f} > \Delta\epsilon^{*}_{f}}$.
        \item Среди оставшихся ${f^{'}_{z}}$ определяется значение, соответствующее минимальному ${N}$. Формируется набор ${\mathbf{s}^{*}}$ соответствующих параметров ${N, M, L}$.
        \item Если нет значений ${f^{'}_{z}}$, удовлетворяющих условию ${\epsilon_{f} < \epsilon^{*}_{f}}$, выбирается набор параметров ${N, M, L}$, соответствующий наименьшему отклонению ${\Delta\epsilon_{f}}$.
    \end{enumerate}
    \item Среди найденных наборов параметров ${\mathbf{s}^{*}}$, соответствующих разным $M$, рассматриваются наборы, удовлетворяющие условию ${\Delta\epsilon_{f} < \Delta\epsilon^{*}_{f}}$ и среди них выбирается набор с минимальным ${N}$.
    \item Если нет наборов параметров ${\mathbf{s}^{*}}$, удовлетворяющих условию ${\Delta\epsilon_{f} < \Delta\epsilon^{*}_{f}}$, выбирается набор параметров, соответствующий наименьшему отклонению ${\Delta\epsilon_{f}}$.
    \item Формируется итоговый набор ${\mathbf{s}^{*} = \left( N, M, L \right)}$.
\end{enumerate}

\textbf{Примечания:}
\begin{enumerate}
    \item Отклонение от заданной частоты ${f_z}$ инжектируемого сигнала допустимо задавать в относительных единицах ${\epsilon^{*}_{f}}$, с последующим пересчетом в абсолютное значение
    \begin{equation*}
        \Delta\epsilon^{*}_{f} = \epsilon^{*}_{f} \cdot f_z.
    \end{equation*}
    В таком случае, на вход алгоритма получается задать общее для всех частот относительное отклонение ${\epsilon^{*}_{f}}$. Абсолютные значения отклонения ${\Delta\epsilon^{*(k)}_{f}}$ сформируются в соответствии частотами ${f^{(k)}_z}$.
    \item Максимальное количество точек ${N_\text{max}}$ определяется размером памяти, доступной в буфере встраиваемой системы управления, зарезервированном под внутрисхемное измерение частотных характеристик.
\end{enumerate}


\textbf{Пример.}

Заданы: ${f_d=100}$~кГц, ${N_\text{min}=10}$, ${N_\text{max}=1214}$.

Определить набор параметров ${N, M, L}$ для измерения на инжектируемой частоте ${f_z=22,1}$~кГц с относительным отклонением до 0,1\% ${\left(\epsilon^{*}_{f}=0,001\right)}$.

Коэффициент соотношения между периодом инжектируемого сигнала и периодом выполнения операций в системе:
\begin{equation*}
    R = \frac{f_{d}}{f_{z}} = \frac{100000}{22100} = 4,524886877828054.
\end{equation*}

Максимально возможное значение коэффициента децимации:
\begin{equation*}
    M_\text{max} = \left\lfloor \frac{R}{2} \right\rfloor = 2.
\end{equation*}

Оптимальный набор параметров: ${N=43}$, ${M=2}$, ${L=19}$.

Соответствующее приближение к начальному значению коэффициента ${R}$:
\begin{equation*}
    R^{'} = \frac{M \cdot N}{\left\lfloor L \right\rfloor} = \frac{2 \cdot 43}{19} = 4,526315789473684.
\end{equation*}

Соответствующее приближение к исходному значению частоты ${f_{z}}$:
\begin{equation*}
    f^{'}_{z} = \frac{L}{M \cdot N}f_d = \frac{19}{2 \cdot 43} = 22093,023255813954~\text{Гц}.
\end{equation*}

Относительное отклонение частоты
\begin{equation*}
    \epsilon_{f} = \frac{\left| f^{'}_{z} - f_{z} \right|}{f_{z}} = 0,0003156897821740197
\end{equation*}
не превышает заданного ограничения ${\epsilon^{*}_{f}=0,001}$.

Таким образом, чтобы воспроизвести частоту инжектируемого сигнала с точностью до 0,1\% достаточно измерить ${N=43}$ точки на ${L=19}$ периодах инжектируемого сигнала. Период выборки при этом должен быть в ${M=2}$ раза больше периода выполнения операций в системе.